% ============================================================================
% MODULE 5: PHENOLOGY & THERMAL TIME
% ============================================================================
% Duration: 60 minutes (Zadok stages, GDD, temperature dependence)
% Learning Goals: Thermal time drives development; Zadok stages; sowing date effects
% Prerequisites: Module 1

\chapter{Module 5: Phenology \& Thermal Time}
\label{ch:module5}

\section{Learning Objectives}

By the end of this module, you will:
\begin{itemize}
  \item Understand thermal time (Growing Degree Days) as driver of crop development
  \item Recognize Zadok stages (0--90+) as standardized phenological scale
  \item Explain why temperature, not calendar, matters for development
  \item Use Zadok stages to identify critical periods (booting, grain fill)
  \item Apply sowing date effects to grain quality and yield
\end{itemize}

\section{Thermal Time: Concept and Calculation}

\subsection{Growing Degree Days (GDD)}

\[
\text{GDD}_{\text{daily}} = \max(0, T_{\text{avg}} - T_{\text{base}})
\]

where:
\begin{itemize}
  \item $T_{\text{avg}}$ = (T$_{\text{max}}$ + T$_{\text{min}}$) / 2
  \item $T_{\text{base}}$ = base temperature (0°C or 3°C depending on crop stage)
  \item Cumulative: $\text{GDD}_{\text{total}} = \sum \text{GDD}_{\text{daily}}$
\end{itemize}

\subsection{Why GDD Instead of Calendar Days?}

\begin{keyconceptbox}
Calendar-based timing fails across climates:
\begin{itemize}
  \item Mediterranean climate: Cool spring (slow development)
  \item Subtropical climate: Warm spring (fast development)
  \item Same calendar date $\ne$ same crop stage across locations
\end{itemize}

GDD standardizes: Same GDD accumulation $\approx$ same physiological development.
\end{keyconceptbox}

\section{Zadok Stages}

\subsection{Standardized Scale (0--90+)}

\begin{table}[H]
\centering
\caption{Zadok Stage Groups and Thermal Time (wheat, subtropical conditions)}
\label{tab:zadok}
\begin{tabularx}{\textwidth}{r|X|r|r}
\toprule
\textbf{Zadok Stage} & \textbf{Description} & \textbf{GDD (cumulative)} & \textbf{Days after sowing} \\
\midrule
0--9  & Germination / Imbibition & 0--50   & 0--10 \\
10--19 & Emergence (seedling leaves) & 50--100  & 10--20 \\
20--29 & Leaf development, tiller initiation & 100--300 & 20--60 \\
30--39 & Stem elongation (growing point differentiates) & 300--600 & 60--120 \\
40--49 & Booting (flag leaf visible, head in boot) & 600--750 & 120--150 \\
50--59 & Heading / head emergence & 750--900 & 150--180 \\
60--69 & Anthesis / pollination & 850--950 & 170--195 \\
70--79 & Grain development (milk to soft dough) & 950--1100 & 195--240 \\
80--89 & Grain ripening (hard dough to hard grain) & 1100--1200 & 240--280 \\
90+   & Physiological maturity / harvest-ready & 1200--1300+ & 280--310 \\
\bottomrule
\end{tabularx}
\end{table}

\begin{tipbox}
\textbf{Check against your simulation:} In your APSIM output, the column \variable{[Wheat].Stage} reports Zadok stage as a decimal (e.g., 45.3). Use this to verify your crop is developing at the expected rate.
\end{tipbox}

\section{Critical Periods}

Not all growth stages are equally sensitive to stress. Yield loss depends strongly on \emph{when} stress occurs:

\begin{center}
\begin{tabular}{l|l|l|l}
\textbf{Stage} & \textbf{Zadok} & \textbf{Stress type} & \textbf{Effect on yield} \\
\hline
Emergence     & 0--10  & Water (dry soil) & Crop fails to establish; zero yield \\
Tillering     & 20--39 & N, water         & Fewer tillers; fewer potential heads \\
Booting       & 40--49 & N, water         & Reduced floret number; irreversible \\
Anthesis      & 55--65 & Water, heat      & Pollen sterility; large grain number loss \\
Grain fill    & 70--85 & Water, heat      & Smaller grain; 30--50\% yield loss if severe \\
Maturity      & 90+    & Rain (excess)    & Grain sprouting in head; quality loss \\
\end{tabular}
\end{center}

\begin{keyconceptbox}
\textbf{Anthesis is the most sensitive single stage.} A hot dry day (T$_{max}$ $>$ 32°C, ESW declining) at Stage 60--65 can cause widespread pollen sterility and reduce final grain number by 20--50\%, regardless of how good the rest of the season is. In APSIM, this is captured through heat and water stress multipliers applied to grain set.
\end{keyconceptbox}

% ---- IMAGE PLACEHOLDER ----
\begin{figure}[H]
  \centering
  \fbox{\parbox{0.85\textwidth}{\centering\vspace{1.5cm}
    \textit{[IMAGE: Bar chart showing relative yield sensitivity to water stress at each Zadok stage.}\\
    \textit{X-axis: Zadok stage groups. Y-axis: Relative yield impact (\%). Anthesis bar is tallest.]}
  \vspace{1.5cm}}}
  \caption{Relative sensitivity of wheat yield to water stress at different Zadok stages. Anthesis and early grain fill are the highest-risk periods.}
  \label{fig:stress-sensitivity}
\end{figure}
% ---- END PLACEHOLDER ----

\section{Sowing Date Effects}

\subsection{Why Sowing Date Changes Everything}

Shifting the sowing date by 4--8 weeks shifts the entire crop calendar. Crucially, it shifts \emph{which temperatures} the crop experiences at each critical stage:

\begin{itemize}
  \item \textbf{Early May sowing} (Dalby): Anthesis falls in August — cool, lower heat stress risk; but risk of frost at booting (July)
  \item \textbf{Mid-June sowing}: Anthesis falls in September — mild temperatures; lowest overall risk; typically optimal for the region
  \item \textbf{Late July sowing}: Anthesis falls in October — warming; higher heat stress risk; shorter grain fill before summer heat
\end{itemize}

% ---- IMAGE PLACEHOLDER ----
\begin{figure}[H]
  \centering
  \fbox{\parbox{0.85\textwidth}{\centering\vspace{1.5cm}
    \textit{[IMAGE: Timeline diagram showing three sowing scenarios. Each is a horizontal bar from Sowing to Maturity.}\\
    \textit{Annotate: anthesis date, frost risk window (Jun--Jul), heat risk window (Oct--Nov) for each scenario.]}
  \vspace{1.5cm}}}
  \caption{Effect of sowing date on crop calendar and risk windows. Shifting sowing by 4 weeks moves anthesis into a different temperature regime.}
  \label{fig:sowing-date-timeline}
\end{figure}
% ---- END PLACEHOLDER ----

\subsection{Exercise C: Sowing Date Factorial}

Three sowing date variants:
\begin{itemize}
  \item Early (May 1--30): Cool spring; long season
  \item Mid (June 15--July 15): Moderate; standard
  \item Late (August 1--31): Warm spring; short season
\end{itemize}

Expected outcomes:
\begin{itemize}
  \item Early sowing: Higher tiller number, grain number; risk of frost at anthesis
  \item Mid sowing: Balanced phenology; maximum grain fill
  \item Late sowing: Short season; less grain per head; risk of terminal drought
\end{itemize}

\section{Temperature Stress}

\subsection{Heat Stress ($T > 30$°C)}

Heat accelerates GDD accumulation (more degree-days per day) but can also damage pollen at anthesis:
\begin{itemize}
  \item Each extra 1°C above T$_\text{base}$ adds 1 extra GDD per day
  \item At T$_{max}$ $>$ 32°C during anthesis, pollen viability drops sharply
  \item Result: Crop matures faster but may set fewer grains
\end{itemize}

\subsection{Cold Stress ($T < T_\text{base}$)}

On days where T$_{avg}$ $<$ T$_\text{base}$ (0--3°C for wheat), GDD does not accumulate:
\begin{itemize}
  \item Development pauses; the crop's timeline extends
  \item If T$_{min}$ $<$ $-$2°C during booting or heading, frost can damage florets irreversibly
  \item In APSIM, frost damage is applied if temperature falls below a crop-specific threshold during vulnerable stages
\end{itemize}

\subsection{Vernalization (Winter Wheat)}

Some wheat varieties require a cold period (vernalization) to trigger flowering. APSIM tracks a vernalization state: the crop will not advance to reproductive stages until sufficient chilling has accumulated. This is primarily relevant for winter wheat in cooler climates; spring wheat varieties (common in subtropical Queensland) have no or minimal vernalization requirement.

\begin{tipbox}
\textbf{Dalby context:} In subtropical Queensland, frost events occur in June--August (particularly in July). An early May sowing means booting occurs in July, creating frost risk at a vulnerable stage. A late July sowing avoids this but shortens the growing season.
\end{tipbox}

\section{Connecting Phenology to Management}

Understanding phenology helps time key management actions:

\begin{center}
\begin{tabular}{l|l|l}
\textbf{Management action} & \textbf{Target Zadok stage} & \textbf{Reason} \\
\hline
Sowing & Stage 0 & Starting point; everything follows from here \\
N fertiliser application & Stage 30--40 & N demand rising; soil moisture often present \\
Late N (topdress) & Stage 40--50 & Just before booting; directly supports grain protein \\
Insecticide/fungicide & Stage 50--65 & Head/flag leaf protection; avoid anthesis disruption \\
Harvest decision & Stage 90--92 & Grain at maximum dry weight; avoid over-drying \\
\end{tabular}
\end{center}

\begin{warningbox}
\textbf{In APSIM, the Manager component applies N at sowing} (Stage 0). In reality, split applications often give better results. For this tutorial, we use a single sowing-time application for simplicity.
\end{warningbox}

\section{Summary}

\begin{keyconceptbox}
\textbf{Module 5 key points:}
\begin{enumerate}
  \item GDD = $\max(0, T_{avg} - T_{base})$ accumulated daily; drives all phenological development
  \item Zadok scale (0--90+) is the universal language for describing wheat stages
  \item Key milestones: Emergence (Stage 10, $\sim$50 GDD), Anthesis (Stage 60, $\sim$900 GDD), Maturity (Stage 90, $\sim$1200 GDD)
  \item Anthesis is the most stress-sensitive stage; N and water stress here cause the largest yield losses
  \item Sowing date shifts the entire crop calendar: early sowing risks frost at booting; late sowing risks heat at anthesis
  \item Time N applications and management actions using Zadok stage, not calendar date
\end{enumerate}
\end{keyconceptbox}

\section{What's Next?}

Module 6 explores long-term climate variability and risk assessment.

% ============================================================================
% END MODULE 5
% ============================================================================
